\documentclass[aps,pra,preprint,a4paper,amsmath,amssymb,longbibliography]{revtex4-2}

\usepackage[T2A]{fontenc}
\usepackage[utf8]{inputenc}
\usepackage[main=russian,english]{babel}
\usepackage{booktabs}
\usepackage{url}
% preprint-стиль REVTeX вставляет перед аннотацией склейку с fil и
% штраф; вместе с babel она уносит аннотацию на следующую страницу.
\makeatletter
\def\frontmatter@abstract@produce{%
  \par
  \addvspace{\frontmatter@preabstractspace}%
  \begingroup
   \dimen@\baselineskip
   \setbox\z@\vtop{\unvcopy\absbox}%
   \advance\dimen@-\ht\z@\advance\dimen@-\prevdepth
   \@ifdim{\dimen@>\z@}{\vskip\dimen@}{}%
  \endgroup
  \begingroup
   \prep@absbox
   \unvbox\absbox
   \post@absbox
  \endgroup
  \@ifx{\@empty\mini@notes}{}{\mini@notes\par}%
  \addvspace\frontmatter@postabstractspace
}%
\makeatother

\newcommand{\zamech}[1]{\par\medskip\noindent\textbf{Замечание.} #1\par\medskip}
\newcommand{\rf}{\mathrm{rf}}
\newcommand{\eq}{\mathrm{eq}}

\begin{document}

\title{Равновесный сдвиг энергии в накопительном кольце\\ во втором порядке продольного движения}
\author{Черновик}
\affiliation{заметка}

\begin{abstract}
Продольные уравнения движения решаются при квадратичной зависимости периода обращения от отклонения энергии и с учётом бетатронного вклада в период. Из периодичности фазы следует точное условие $\langle\tau\rangle=0$, а метод Боголюбова--Митропольского даёт равновесный сдвиг $\delta_\eq=-(\eta_1\delta_m^2/2+\Lambda)/\eta_0$, вторую гармонику, сдвиг синхротронной частоты и поправку конечной фазы. Для магнитного кольца выведены $\eta_0$, $\eta_1$ и $\Lambda$; для электростатического --- те же коэффициенты с учётом изменения скорости в потенциале дефлектора, с примерами гладких колец с цилиндрическим и сферическим дефлекторами. Решение проверено интегрированием полной системы.
\end{abstract}

\maketitle

\section{Постановка}
\label{sec:setup}

Частица в накопителе совершает синхротронные колебания около синхронной частицы. Если период обращения зависит от отклонения энергии нелинейно или от амплитуды бетатронных колебаний, среднее отклонение энергии за синхротронный период смещается от нуля. Это среднее --- равновесный сдвиг $\delta_\eq$ --- и находится ниже. Схема такая: сначала строится функция $\tau$, относительное опоздание частицы за оборот (разделы~\ref{sec:mag} и~\ref{sec:beta}); затем один раз решаются продольные уравнения с произвольной $\tau$ (раздел~\ref{sec:sol}); электростатическое кольцо меняет только коэффициенты $\tau$ (раздел~\ref{sec:es}).

\subsection{Переменная $\delta$}

Индекс $s$ отмечает синхронную частицу: импульс $p_s$, скорость $v_s=\beta_s c$, лоренц-фактор $\gamma_s$, энергия $E_s=\gamma_s mc^2$, длина орбиты $C_s$, период $T_s=C_s/v_s$. Функции Твисса $\beta_{x,y}$, $\alpha_{x,y}$, $\gamma_{x,y}$ всегда пишутся с индексом плоскости.

Между ускоряющими зазорами поле статическое, и сохраняется полная энергия $E=\gamma mc^2+q\Phi$. Отклонение $\delta$ определяем через неё: $p_s(1+\delta)$ есть импульс, который частица с полной энергией $E$ имела бы там, где $\Phi=0$,
\begin{equation}
\label{eq:set-delta}
E=\sqrt{p_s^2(1+\delta)^2c^2+m^2c^4}.
\end{equation}
В магнитном кольце $\Phi\equiv 0$, и $\delta=(p-p_s)/p_s$ --- обычное отклонение импульса. В электростатическом кольце локальный импульс в дефлекторе меняется вдоль траектории, а $\delta$ между зазорами постоянно и меняется только в зазоре. Поэтому именно $\delta$ --- медленная переменная продольного движения, и равновесный сдвиг есть её среднее по синхротронному колебанию:
\begin{equation}
\label{eq:set-deq}
\delta_\eq=\langle\delta\rangle.
\end{equation}

Длина замкнутой орбиты записывается рядом с коэффициентами $\alpha_0$, $\alpha_1$ (коэффициенты ряда, не производные):
\begin{equation}
\label{eq:set-C}
\frac{C}{C_s}=1+\alpha_0\delta+\alpha_1\delta^2+O(\delta^3).
\end{equation}

\subsection{Опоздание за оборот}

Относительное опоздание
\begin{equation}
\label{eq:set-tau}
\tau=\frac{T-T_s}{T_s}=\eta_0\delta+\eta_1\delta^2+\Lambda
\end{equation}
есть ряд по $\delta$ плюс вклад $\Lambda$ бетатронных колебаний, усреднённый по быстрой бетатронной фазе. Делитель --- синхронный период: за оборот фаза частицы относительно поля смещается ровно на $\omega_\rf(T-T_s)$, и в фазовое уравнение входит именно $(T-T_s)/T_s$.

\zamech{Если коэффициент скольжения определять через $(T-T_s)/T$ или через частоту обращения, $\eta=-\Delta\omega/(\omega_0\delta)$ (формула (18.20) в~\cite{Ng2006}), то во втором порядке вместо $\eta_1$ получается $\eta_1-\eta_0^2$; определение через производную, $\eta=-\omega_0^{-1}\mathrm{d}\omega/\mathrm{d}\delta$ (18.21), даёт $2(\eta_1-\eta_0^2)$. Линейный коэффициент у всех определений общий, но в фазовое уравнение входит только $\eta_1$ из~\eqref{eq:set-tau}.}

\subsection{Уравнения движения}

Фаза $\varphi$ --- отставание от синхронной частицы в радианах ВЧ-поля, она растёт при опоздании. Напряжение зазора $V\sin(\varphi_s+\varphi)$, $\omega_\rf=2\pi h/T_s$. Без ускорения $\varphi_s=0$ или $\pi$; устойчива точка
\begin{equation}
\label{eq:set-phis}
\cos\varphi_s=-\operatorname{sgn}\eta_0 .
\end{equation}
За зазор полная энергия меняется на $eV\cos\varphi_s\sin\varphi$, а $\mathrm{d}E=\beta_s^2E_s\,\mathrm{d}\delta$. Сглаживание по обороту даёт
\begin{equation}
\label{eq:set-system}
\dot\varphi=\omega_\rf\,\tau,
\qquad
\dot\delta=K\cos\varphi_s\sin\varphi,
\qquad
K=\frac{eV\omega_\rf}{2\pi h\,\beta_s^2E_s}.
\end{equation}
Точка обозначает производную по времени синхронной частицы. Линеаризация около $\varphi=\delta=0$ даёт колебания с частотой
\begin{equation}
\label{eq:set-ws}
\omega_s^2=K\omega_\rf|\eta_0|
\end{equation}
именно при выборе~\eqref{eq:set-phis}.

\section{Опоздание в магнитном кольце}
\label{sec:mag}

Магнитная сила перпендикулярна скорости, поэтому между зазорами модуль скорости постоянен вдоль всей траектории. Период $T=\oint\mathrm{d}l/v=C/v$, и
\begin{equation}
\label{eq:mag-Tratio}
\frac{T}{T_s}=\frac{C}{C_s}\,\frac{v_s}{v}
\end{equation}
точно. Скорость задаётся импульсом: из $\beta\gamma=\beta_s\gamma_s(1+\delta)$ и $\gamma^2=1+(\beta\gamma)^2$
\begin{equation}
\label{eq:mag-v}
\frac{v_s}{v}=\frac{\gamma(\delta)}{\gamma_s(1+\delta)}
=1-\frac{\delta}{\gamma_s^2}+\Biggl(\frac{1}{\gamma_s^4}+\frac{3\beta_s^2}{2\gamma_s^2}\Biggr)\delta^2+O(\delta^3).
\end{equation}
Слагаемое $1/\gamma_s^4$ --- квадрат линейного члена, $3\beta_s^2/(2\gamma_s^2)$ --- вторая производная скорости по импульсу; оба одного порядка. Произведение~\eqref{eq:mag-v} и~\eqref{eq:set-C} даёт коэффициенты~\eqref{eq:set-tau}:
\begin{equation}
\label{eq:mag-eta}
\eta_0=\alpha_0-\frac{1}{\gamma_s^2},
\qquad
\eta_1=\alpha_1-\frac{\alpha_0}{\gamma_s^2}+\frac{1}{\gamma_s^4}+\frac{3\beta_s^2}{2\gamma_s^2}.
\end{equation}

\section{Бетатронный вклад}
\label{sec:beta}

При бетатронных колебаниях частица проходит не замкнутую орбиту, а траекторию длины $L=C(\delta)+\Delta L_\beta$. В магнитном кольце скорость вдоль этой траектории та же, что на замкнутой орбите, поэтому~\eqref{eq:mag-Tratio} сохраняется с заменой $C\to L$:
\begin{equation}
\label{eq:beta-tau}
\tau=\eta_0\delta+\eta_1\delta^2+\frac{\Delta L_\beta}{C_s}\Bigl(1+O(\delta)\Bigr).
\end{equation}
Отношение времён равно отношению длин именно потому, что обе длины проходятся с одной скоростью; в электростатическом кольце это не так (раздел~\ref{sec:es}).

Синхротронное движение медленное и видит среднее по бетатронной фазе, $\Lambda=\langle\Delta L_\beta\rangle/C_s$. Элемент пути до второго порядка
\begin{equation}
\label{eq:beta-dl}
\mathrm{d}l=\Biggl(1+\frac{x}{\rho}+\frac{x'^2+y'^2}{2}\Biggr)\mathrm{d}s
\end{equation}
($(x/\rho)^2$ сокращается при разложении корня). Смещение $x=D\delta+x_\beta+x_{\mathrm{nl}}$: дисперсионная часть уже учтена в $\alpha_0$, $\alpha_1$; линейное колебание $x_\beta=\sqrt{2J_x\beta_x}\cos\psi_x$ в среднем даёт нуль в $x/\rho$; $x_{\mathrm{nl}}$ --- сдвиг замкнутой орбиты, квадратичный по амплитуде (от секступолей, кривизны, краевых полей). Среднее квадрата угла $\langle x_\beta'^2\rangle=J_x\gamma_x$, и
\begin{equation}
\label{eq:beta-Lambda}
\Lambda=\frac{1}{C_s}\Biggl(\frac{J_x}{2}\oint\gamma_x\,\mathrm{d}s+\frac{J_y}{2}\oint\gamma_y\,\mathrm{d}s+\Biggl\langle\oint\frac{x_{\mathrm{nl}}}{\rho}\,\mathrm{d}s\Biggr\rangle\Biggr).
\end{equation}
В гладком приближении ($\alpha_{x,y}=0$, $\gamma_x=1/\beta_x$, $\oint\mathrm{d}s/\beta_x=2\pi\nu_x$) и без $x_{\mathrm{nl}}$
\begin{equation}
\label{eq:beta-smooth}
\Lambda=\frac{\pi}{C_s}\bigl(J_x\nu_x+J_y\nu_y\bigr)=\frac{\pi}{2C_s}\bigl(\varepsilon_x\nu_x+\varepsilon_y\nu_y\bigr),
\end{equation}
где $\varepsilon=2J$ --- инвариант Куранта--Снайдера одной частицы.

\zamech{С учётом $x_{\mathrm{nl}}$ та же величина выражается через хроматичности $\xi_{x,y}=\mathrm{d}\nu_{x,y}/\mathrm{d}\delta$: $\Lambda=-(2\pi/C_s)(J_x\xi_x+J_y\xi_y)$~\cite{Shoji2005}.}

\section{Решение продольных уравнений}
\label{sec:sol}

Решаем систему~\eqref{eq:set-system} с $\tau=\eta_0\delta+\eta_1\delta^2+\Lambda$, считая $\eta_1\delta$ и $\Lambda$ малыми по сравнению с $\eta_0\delta$.

\subsection{Точное условие на среднее}

На замкнутой синхротронной траектории фаза через период возвращается к прежнему значению, поэтому $\langle\dot\varphi\rangle=0$, и первое уравнение~\eqref{eq:set-system} даёт
\begin{equation}
\label{eq:sol-mean}
\langle\tau\rangle=0
\quad\Longrightarrow\quad
\delta_\eq=-\frac{\eta_1\langle\delta^2\rangle+\Lambda}{\eta_0}.
\end{equation}
Среднее берётся по времени за синхротронный период. Соотношение точное: оно не зависит ни от формы напряжения, ни от приближений во втором уравнении, ни от знака $\eta_0$. Всё остальное решение нужно лишь для того, чтобы выразить $\langle\delta^2\rangle$ через амплитуду колебаний.

Смысл знака: при $\eta_1>0$ член $\eta_1\delta^2$ удлиняет средний период при любом знаке $\delta$; чтобы среднее опоздание было нулевым, линейный член должен его скомпенсировать, $\eta_0\delta_\eq<0$. Ниже перехода ($\eta_0<0$) равновесный импульс поднимается.

\subsection{Уравнение для $\delta$}

Дифференцирование второго уравнения~\eqref{eq:set-system} даёт $\ddot\delta=K\omega_\rf\cos\varphi_s\cos\varphi\cdot\tau$. С~\eqref{eq:set-phis} и~\eqref{eq:set-ws} $K\omega_\rf\cos\varphi_s=-\omega_s^2/\eta_0$, и
\begin{equation}
\label{eq:sol-exact}
\ddot\delta=-\frac{\omega_s^2}{\eta_0}\cos\varphi\,\bigl(\eta_0\delta+\eta_1\delta^2+\Lambda\bigr).
\end{equation}
При малом размахе фазы $\cos\varphi\approx 1$:
\begin{equation}
\label{eq:sol-reduced}
\ddot\delta+\omega_s^2\delta=c_0+c_2\delta^2,
\qquad
c_0=-\frac{\omega_s^2\Lambda}{\eta_0},
\quad
c_2=-\frac{\omega_s^2\eta_1}{\eta_0}.
\end{equation}

\subsection{Решение до второго порядка}

Метод Боголюбова--Митропольского (приложение~\ref{app:bm}) для~\eqref{eq:sol-reduced} даёт
\begin{gather}
\label{eq:sol-delta}
\delta=\delta_m\cos\psi+\delta_\eq+\frac{\eta_1\delta_m^2}{6\eta_0}\cos 2\psi,
\qquad
\dot\psi=\Omega,
\\
\label{eq:sol-deq}
\delta_\eq=-\frac{1}{\eta_0}\Biggl(\frac{\eta_1\delta_m^2}{2}+\Lambda\Biggr),
\\
\label{eq:sol-omega}
\frac{\Omega}{\omega_s}=1-\frac{\eta_1\Lambda}{\eta_0^2}-\frac{5}{12}\Biggl(\frac{\eta_1\delta_m}{\eta_0}\Biggr)^2.
\end{gather}
Здесь $\delta_m$ --- амплитуда первой гармоники. Формула~\eqref{eq:sol-deq} совпадает с~\eqref{eq:sol-mean} при $\langle\delta^2\rangle=\delta_m^2/2$.

\textbf{Вторая гармоника} в $\delta_\eq$ на порядке $\delta_m^2$ не входит. Она и сам сдвиг поправляют $\langle\delta^2\rangle$ на $\delta_\eq^2+\tfrac12(\eta_1\delta_m^2/6\eta_0)^2$, что при $\Lambda=0$ даёт следующий член
\begin{equation}
\label{eq:sol-a4}
\delta_\eq=-\frac{\eta_1\delta_m^2}{2\eta_0}-\frac{19}{72}\,\frac{\eta_1^3\delta_m^4}{\eta_0^3}.
\end{equation}

\textbf{Сдвиг частоты}~\eqref{eq:sol-omega} состоит из двух частей. Член $-\eta_1\Lambda/\eta_0^2$ виден уже при нулевой амплитуде: рабочая точка сдвинута в $\delta_\eq=-\Lambda/\eta_0$, где наклон $\tau$ равен $\eta_0+2\eta_1\delta_\eq$. Член $-5/12$ --- эффект конечной амплитуды: $1/2$ от постоянного смещения и $-1/12$ от второй гармоники. Он отрицателен при любом знаке $\eta_1/\eta_0$.

\subsection{Поправка конечной фазы}

Удержим в~\eqref{eq:sol-exact} $\cos\varphi\approx 1-\varphi^2/2$. На линейной орбите $\delta=\delta_m\cos\psi$ интегрирование $\dot\varphi=\omega_\rf\eta_0\delta$ даёт $\varphi=\pm\varphi_m\sin\psi$ с
\begin{equation}
\label{eq:sol-phim}
\varphi_m=\frac{\omega_\rf|\eta_0|}{\omega_s}\,\delta_m .
\end{equation}
Добавка к правой части $\omega_s^2(\varphi^2/2)\delta=\omega_s^2(\varphi_m^2\delta_m/8)(\cos\psi-\cos 3\psi)$ не имеет постоянной составляющей, поэтому среднее не сдвигает; её первая гармоника сдвигает частоту на
\begin{equation}
\label{eq:sol-dphi}
\frac{\Delta\Omega}{\omega_s}=-\frac{\varphi_m^2}{16}.
\end{equation}
При $\eta_1=\Lambda=0$ система~\eqref{eq:set-system} --- маятник $\ddot\varphi+\omega_s^2\sin\varphi=0$, и~\eqref{eq:sol-dphi} --- первый член его известного разложения $\Omega/\omega_s=1-\varphi_m^2/16+\varphi_m^4/3072-\dots$; $\langle\delta\rangle=0$ при любой амплитуде.

В $\delta_\eq$ поправка фазы входит только через $\langle\delta^2\rangle$: третья гармоника амплитуды $\varphi_m^2\delta_m/64$ меняет его на величину порядка $\varphi_m^4\delta_m^2$. Итого:
\begin{equation}
\label{eq:sol-omega-sum}
\frac{\Omega}{\omega_s}=1-\frac{\eta_1\Lambda}{\eta_0^2}-\frac{5}{12}\Biggl(\frac{\eta_1\delta_m}{\eta_0}\Biggr)^2-\frac{\varphi_m^2}{16},
\end{equation}
а сдвиг $\delta_\eq$ даётся~\eqref{eq:sol-deq} с поправками порядка $\delta_m^4$ и $\varphi_m^4\delta_m^2$.

\subsection{Применимость}

Разложение требует
\begin{equation}
\label{eq:sol-valid}
\Biggl|\frac{\eta_1\delta_m}{\eta_0}\Biggr|\ll 1,
\qquad
\Biggl|\frac{\Lambda}{\eta_0\delta_m}\Biggr|\ll 1,
\qquad
\varphi_m\ll 1.
\end{equation}
Вблизи перехода $\eta_0\to 0$ первое условие нарушается при любой амплитуде; условие~\eqref{eq:sol-mean} тогда превращается в $\eta_1\langle\delta^2\rangle+\Lambda=0$, и для сдвига нужен следующий член ряда $\tau$.

Численная проверка решения на полной системе~\eqref{eq:set-system} приведена в приложении~\ref{app:num}.

\section{Электростатическое кольцо}
\label{sec:es}

Зазоры стоят там, где $\Phi=0$, поэтому зазор меняет полную энергию, то есть $\delta$ из~\eqref{eq:set-delta}, на ту же величину, что в магнитном кольце. Уравнения~\eqref{eq:set-system} и всё решение раздела~\ref{sec:sol} не меняются. Меняется только функция $\tau$: в дефлекторе сила имеет составляющую вдоль скорости, и на смещённой траектории частица движется с другой скоростью. Поэтому опоздание за оборот складывается из удлинения пути и вклада изменённой скорости, а $v$ в $T=\oint\mathrm{d}l/v$ из-под интеграла не выносится.

\subsection{Потенциал и скорость}

Потенциал дефлектора с опорной дугой радиуса $\rho$ до второго порядка, согласованный с уравнением Лапласа и с условием $qE_0=p_sv_s/\rho$:
\begin{equation}
\label{eq:es-phi}
\psi\equiv\frac{q\Phi}{p_sv_s}=\frac{x}{\rho}-\frac{n}{2}\Bigl(\frac{x}{\rho}\Bigr)^2+\frac{n-1}{2}\Bigl(\frac{y}{\rho}\Bigr)^2.
\end{equation}
Цилиндрический дефлектор ($E\propto1/r$) соответствует $n=1$, сферический ($E\propto1/R^2$) --- $n=2$. Локальный лоренц-фактор $\gamma=\gamma_E(\delta)-\beta_s^2\gamma_s\psi$, и с $\psi=\psi_1+\psi_2$ (линейная и квадратичная по отклонению части)
\begin{equation}
\label{eq:es-speed}
\frac{v_s}{v}=1-\frac{\delta-\psi}{\gamma_s^2}-\frac{\delta^2}{2\gamma_s^4}+\frac{3}{2\gamma_s^2}(\delta-\psi_1)^2.
\end{equation}
При $\psi=0$ это~\eqref{eq:mag-v}. Опоздание
\begin{equation}
\label{eq:es-tau}
\tau=\frac{1}{C_s}\oint\Biggl(\frac{v_s}{v}\,\frac{\mathrm{d}l}{\mathrm{d}s}-1\Biggr)\mathrm{d}s
\end{equation}
с элементом пути~\eqref{eq:beta-dl}.

\subsection{Коэффициенты $\eta_0$, $\eta_1$}

На замкнутой орбите $x=D\delta+D_2\delta^2$. В первом порядке подынтегральное выражение~\eqref{eq:es-tau} равно $-\delta/\gamma_s^2+x/\rho+(x/\rho_E)/\gamma_s^2$, где $1/\rho_E$ --- кривизна, создаваемая электрическим полем. Отсюда
\begin{equation}
\label{eq:es-eta0}
\eta_0=\alpha_0+\frac{\alpha_0^E}{\gamma_s^2}-\frac{1}{\gamma_s^2},
\qquad
\alpha_0^E=\frac{1}{C_s}\oint\frac{D}{\rho_E}\,\mathrm{d}s .
\end{equation}
Смещение наружу и удлиняет путь, и через потенциал уменьшает скорость; второе даёт $\alpha_0^E/\gamma_s^2$. В чисто электрическом кольце $\alpha_0^E=\alpha_0$, $\eta_0=(1+1/\gamma_s^2)\alpha_0-1/\gamma_s^2$, и переход лежит при $\gamma_s^2=1/\alpha_0-1$ (только если $\alpha_0<1$). В фильтре Вина суммарная кривизна орбиты нулевая, но электрическое поле есть: участок не даёт вклада в $\alpha_0$ и даёт вклад в $\alpha_0^E$.

Во втором порядке для кольца с электрическим изгибом ($\rho_E=\rho$)
\begin{equation}
\label{eq:es-eta1}
\eta_1=\eta_1^{\mathrm{mag}}+\frac{1}{\gamma_s^2}\Biggl(\Biggl\langle\frac{D_2}{\rho}\Biggr\rangle+\frac{5-n}{2}\Biggl\langle\Bigl(\frac{D}{\rho}\Bigr)^2\Biggr\rangle-3\alpha_0\Biggr),
\end{equation}
где $\eta_1^{\mathrm{mag}}$ --- формула~\eqref{eq:mag-eta} на тех же $\alpha_0$, $\alpha_1$, а $\langle f\rangle=C_s^{-1}\oint f\,\mathrm{d}s$. Добавка не выражается через $\alpha_0$, $\alpha_1$: в неё входит только $D_2$ (а не сумма $D_2/\rho+D'^2/2$, составляющая $\alpha_1$) и второй момент дисперсии. В смешанном кольце добавка интегрируется по электрическим дугам, а $-3\alpha_0$ заменяется на $-3\alpha_0^E$.

\subsection{Бетатронный вклад}

При $\delta=0$ в~\eqref{eq:es-speed} остаётся $\psi$ на бетатронной траектории $x=x_\beta+x_{\mathrm{nl}}$. Среднее по бетатронной фазе даёт
\begin{equation}
\label{eq:es-lambda}
\Lambda=\frac{1}{C_s}\oint\Biggl[\Bigl(1+\frac{1}{\gamma_s^2}\Bigr)\frac{\langle x_{\mathrm{nl}}\rangle}{\rho}+\frac{\langle x'^2+y'^2\rangle}{2}+\frac{5-n}{2\gamma_s^2}\frac{\langle x_\beta^2\rangle}{\rho^2}+\frac{n-1}{2\gamma_s^2}\frac{\langle y^2\rangle}{\rho^2}\Biggr]\mathrm{d}s.
\end{equation}
Вне дефлекторов остаётся только угловой член, как в~\eqref{eq:beta-Lambda}. Внутри смещение $x_{\mathrm{nl}}$ входит с весом $1+1/\gamma_s^2$ (путь плюс замедление), а квадраты смещений, которых в длине пути нет, входят через кривизну потенциала. Поэтому $\Lambda$ в электростатике не пропорционально удлинению пути $\Delta L_\beta/C_s$.

\subsection{Гладкие кольца}

Кольцо из одного дефлектора без прямых промежутков позволяет найти всё точно из баланса силы и энергии на круговых орбитах. Результаты (выраженные через $X=\langle(x_\beta/\rho)^2\rangle$):

\begin{table}[htb]
\caption{Гладкое кольцо из одного дефлектора.}
\label{tab:es-smooth}
\begin{ruledtabular}
\begin{tabular}{lcc}
& цилиндрический ($n=1$) & сферический ($n=2$) \\
\hline
$\alpha_0$ & $1$ & $\gamma_s^2+1$ \\
$\eta_0$ & $1$ & $\gamma_s^2+2$ \\
$\eta_1$ & $\dfrac{\gamma_s^2+1}{2\gamma_s^2}$ & $\dfrac{9\gamma_s^4+7\gamma_s^2+2}{2\gamma_s^2}$ \\
$\nu_x$ & $\sqrt{2-\beta_s^2}$ & $1/\gamma_s$ \\
$\langle x_{\mathrm{nl}}\rangle/\rho$ & $-\dfrac{\gamma_s^2+5}{2(\gamma_s^2+1)}X$ & $-X$ \\
$\Lambda$ & $0$ & $-\beta_s^2X$ \\
$\Delta L_\beta/C_s$ & $\dfrac{1-3\gamma_s^2}{2\gamma_s^2(\gamma_s^2+1)}X$ & $\Bigl(\dfrac{1}{2\gamma_s^2}-1\Bigr)X$ \\
\end{tabular}
\end{ruledtabular}
\end{table}

В цилиндрическом дефлекторе $pv=\mathrm{const}$ на всех круговых орбитах, скорость на них равна $v_s$, и весь избыток энергии уходит в потенциал: $r/\rho=\exp(\delta+\delta^2/2\gamma_s^2+\dots)$, $\tau=r/\rho-1$. Бетатронное колебание удлиняет путь, но в той же мере меняет среднюю скорость, и $\Lambda=0$. В сферическом дефлекторе $\Lambda\neq0$ и отличается от удлинения пути. Вертикальное колебание в нём --- наклон плоскости круговой орбиты, и его вклад в период равен нулю точно.

Подстановка в~\eqref{eq:sol-deq}:
\begin{equation}
\label{eq:es-deq}
\delta_\eq^{\mathrm{cyl}}=-\frac{\gamma_s^2+1}{4\gamma_s^2}\,\delta_m^2,
\qquad
\delta_\eq^{\mathrm{sph}}=-\frac{1}{\gamma_s^2+2}\Biggl(\frac{9\gamma_s^4+7\gamma_s^2+2}{4\gamma_s^2}\,\delta_m^2-\beta_s^2X\Biggr).
\end{equation}

\subsection{Локальная энергия}

Сдвиг $\delta_\eq$ относится к полной энергии. Кинетическая энергия внутри дефлектора, которую видит, например, спин, --- другая величина. На дисперсионной орбите чисто электрического кольца её среднее по обороту в единицах $\delta$ равно $(1-\alpha_0)\delta+O(\delta^2)$: доля $\alpha_0$ отклонения энергии сидит в потенциале. В гладком кольце с цилиндрическим дефлектором $\alpha_0=1$, и кинетическая энергия на круговой орбите не меняется вовсе: синхротронное колебание и сдвиг~\eqref{eq:es-deq} меняют радиус и полную энергию, но не локальный $\gamma$. Бетатронные колебания, напротив, сдвигают среднюю кинетическую энергию и при $\Lambda=0$: в том же кольце $\langle\gamma-\gamma_s\rangle/(\beta_s^2\gamma_s)=X(\gamma_s^2+3)/(\gamma_s^2+1)$.


\appendix
\section{Метод Боголюбова--Митропольского}
\label{app:bm}

Здесь собрана только та часть метода~\cite{Bogoliubov1961}, которой пользуется раздел~\ref{sec:sol}. Решается
\begin{equation}
\label{eq:bm-eq}
\ddot x+\omega^2 x=\varepsilon f(x,\dot x)
\end{equation}
с малым параметром $\varepsilon$. Ищем
\begin{align}
x&=a\cos\psi+\varepsilon u_1(a,\psi)+\varepsilon^2 u_2(a,\psi)+\cdots,
\label{eq:bm-ansatz}\\
\dot a&=\varepsilon A_1(a)+\varepsilon^2 A_2(a)+\cdots,
\label{eq:bm-adot}\\
\dot\psi&=\omega+\varepsilon B_1(a)+\varepsilon^2 B_2(a)+\cdots.
\label{eq:bm-psidot}
\end{align}
Функции $u_k$ периодичны по $\psi$ с периодом $2\pi$ и не содержат первой гармоники: она отнесена к амплитуде $a$ и фазе $\psi$. Если оставить $\cos\psi$ или $\sin\psi$ внутри $u_k$, оператор $\partial_{\psi\psi}+1$ на них обращается в нуль, и поправка растёт линейно по времени.

В применении к продольному движению правая часть уже мала, и $\varepsilon$ служит только меткой порядка. В окончательных формулах $\varepsilon=1$, а малость сидит в коэффициентах $c_0$, $c_2$: поправка первого порядка линейна по ним, второго --- квадратична. Подстановка ряда Линдштедта $x=a\cos\Omega t+\cdots$ даёт те же коэффициенты; это независимая проверка выкладок ниже.

\subsection{Первый порядок}
\label{sec:bm-first}

На невозмущённой орбите $x=a\cos\psi$, $\dot x=-a\omega\sin\psi$. Подстановка~\eqref{eq:bm-ansatz}--\eqref{eq:bm-psidot} и сбор членов порядка $\varepsilon$ дают
\begin{equation}
\label{eq:bm-order1}
\omega^2(\partial_{\psi\psi}u_1+u_1)
-2\omega A_1\sin\psi
-2a\omega B_1\cos\psi
=f_0(\psi),
\end{equation}
где $f_0(\psi)=f(a\cos\psi,-a\omega\sin\psi)$. Множитель $2$ возникает потому, что производная по времени действует и на медленную амплитуду, и на фазу: поправка скорости $A_1\cos\psi-a B_1\sin\psi$ дифференцируется ещё раз со скоростью $\omega$.

Ряд Фурье правой части записываем так:
\begin{equation}
\label{eq:bm-fourier}
f_0=g_0+\sum_{n\ge 1}(g_n\cos n\psi+h_n\sin n\psi),
\end{equation}
\begin{equation}
\label{eq:bm-fourier-coeff}
\begin{aligned}
g_0&=\frac{1}{2\pi}\int_0^{2\pi}f_0\,\mathrm{d}\psi,\\
g_n&=\frac{1}{\pi}\int_0^{2\pi}f_0\cos n\psi\,\mathrm{d}\psi,
\qquad
h_n=\frac{1}{\pi}\int_0^{2\pi}f_0\sin n\psi\,\mathrm{d}\psi.
\end{aligned}
\end{equation}
Нормировка $g_0$ отличается от $g_n$ вдвое: в разложении постоянный член стоит без дополнительного множителя $1/2$.

Коэффициент при $\cos\psi$ в~\eqref{eq:bm-order1} не может быть уравновешен функцией $u_1$, и то же верно для $\sin\psi$. Эти два члена переносятся в эволюцию амплитуды и фазы:
\begin{equation}
\label{eq:bm-ab}
A_1=-\frac{h_1}{2\omega},
\qquad
B_1=-\frac{g_1}{2a\omega}.
\end{equation}
Остаток ряда определяет поправку формы
\begin{equation}
\label{eq:bm-u}
u_1=\frac{g_0}{\omega^2}
+\sum_{n\ge 2}\frac{g_n\cos n\psi+h_n\sin n\psi}{\omega^2(1-n^2)}.
\end{equation}
Знаменатель $1-n^2$ есть собственное число $\partial_{\psi\psi}+1$ на гармонике $n$.

\subsection{Квадратичная правая часть}
\label{sec:bm-quad}

Для
\begin{equation}
\label{eq:bm-quad-f}
f=c_0+c_2 x^2
\end{equation}
на орбите $x=a\cos\psi$ тождество $\cos^2\psi=(1+\cos 2\psi)/2$ даёт
\begin{equation}
\label{eq:bm-f0}
f_0=c_0+\frac{c_2 a^2}{2}+\frac{c_2 a^2}{2}\cos 2\psi.
\end{equation}
Ненулевые коэффициенты --- $g_0=c_0+c_2 a^2/2$ и $g_2=c_2 a^2/2$. Первой гармоники нет, поэтому
\begin{equation}
\label{eq:bm-ab1-zero}
A_1=0,
\qquad
B_1=0,
\end{equation}
а~\eqref{eq:bm-u} превращается в
\begin{equation}
\label{eq:bm-u1}
u_1
=\frac{c_0+c_2 a^2/2}{\omega^2}
-\frac{c_2 a^2}{6\omega^2}\cos 2\psi.
\end{equation}
Второе слагаемое получается так: $(1-4)U_2=g_2/\omega^2$, откуда $U_2=-g_2/(3\omega^2)=-c_2 a^2/(6\omega^2)$. Постоянная часть $u_1$ есть среднее $\langle x\rangle$ на этом порядке. Сдвига частоты на первом порядке нет: квадрат $x^2$ не проецируется на $\cos\psi$.

\subsection{Второй порядок}
\label{sec:bm-second}

Поскольку $A_1=B_1=0$, амплитуда и фаза начинают меняться с $A_2$ и $B_2$. Подстановка $x=a\cos\psi+u_1+u_2$ в $c_0+c_2 x^2$ добавляет к правой части член $2c_2(a\cos\psi)\,u_1$: произведение $c_2 u_1^2$ имеет следующий порядок малости. Уравнение того же вида, что~\eqref{eq:bm-order1}, читается
\begin{equation}
\label{eq:bm-order2}
\omega^2(\partial_{\psi\psi}u_2+u_2)
-2\omega A_2\sin\psi
-2a\omega B_2\cos\psi
=2c_2 a\cos\psi\cdot u_1.
\end{equation}
Запись $u_1=U_0+U_2\cos 2\psi$ и тождество $\cos\psi\cos 2\psi=(\cos 3\psi+\cos\psi)/2$ дают
\begin{equation}
\label{eq:bm-f1}
2c_2 a\cos\psi\cdot u_1
=2c_2 a\Bigl(U_0+\frac{U_2}{2}\Bigr)\cos\psi
+c_2 a U_2\cos 3\psi.
\end{equation}
Синусов нет, поэтому $A_2=0$: первая гармоника решения не затухает и не растёт. Косинус первой гармоники снимается сдвигом частоты по правилу~\eqref{eq:bm-ab},
\begin{equation}
\label{eq:bm-b2-step}
B_2=-\frac{c_2}{\omega}\Bigl(U_0+\frac{U_2}{2}\Bigr).
\end{equation}
Подстановка $U_0$ и $U_2$ из~\eqref{eq:bm-u1} собирает
\begin{equation}
\label{eq:bm-split}
\begin{aligned}
U_0+\frac{U_2}{2}
&=\frac{c_0}{\omega^2}+\frac{c_2 a^2}{\omega^2}\Bigl(\frac12-\frac{1}{12}\Bigr)
=\frac{c_0}{\omega^2}+\frac{5}{12}\frac{c_2 a^2}{\omega^2}.
\end{aligned}
\end{equation}
Слагаемое $1/2$ --- постоянное смещение, порождённое средним $x^2$. Слагаемое $-1/12$ --- вторая гармоника: $U_2/2=-c_2 a^2/(12\omega^2)$. Отсюда
\begin{equation}
\label{eq:bm-b2}
B_2=-\frac{c_2 c_0}{\omega^3}-\frac{5}{12}\frac{c_2^2 a^2}{\omega^3},
\end{equation}
и частота колебания есть $\dot\psi=\omega+B_2$.

Постоянного члена в~\eqref{eq:bm-f1} нет, поэтому постоянная часть $u_2$ равна нулю. Среднее $\langle x\rangle$ через второй порядок остаётся постоянной частью $u_1$. Третья гармоника из~\eqref{eq:bm-f1} определяет форму, но не частоту и не среднее:
\begin{equation}
\label{eq:bm-v3}
u_2=\frac{c_2^2 a^3}{48\omega^4}\cos 3\psi.
\end{equation}
Её квадрат входит в $\langle x^2\rangle$ только на следующем порядке.

Разложение пригодно, пока поправка мала по сравнению с линейной возвращающей силой,
\begin{equation}
\label{eq:bm-valid}
\biggl|\frac{c_2 a}{\omega^2}\biggr|\ll 1,
\qquad
\biggl|\frac{c_0}{\omega^2 a}\biggr|\ll 1.
\end{equation}
Второе условие относит постоянную силу к масштабу $\omega^2 a$.

Подстановка $\omega=\omega_s$, $a=\delta_m$ и коэффициентов~\eqref{eq:sol-reduced} в~\eqref{eq:bm-u1} и~\eqref{eq:bm-b2} даёт~\eqref{eq:sol-delta}--\eqref{eq:sol-omega}.

\section{Численная проверка}
\label{app:num}

\subsection{Продольное движение}

С безразмерным временем $\theta=\omega_s t$ и $\mu=\omega_\rf/\omega_s$ система~\eqref{eq:set-system} принимает вид
\begin{equation}
\label{eq:num-sys}
\frac{\mathrm{d}\varphi}{\mathrm{d}\theta}=\mu\bigl(\eta_0\delta+\eta_1\delta^2+\Lambda\bigr),
\qquad
\frac{\mathrm{d}\delta}{\mathrm{d}\theta}=-\frac{\sin\varphi}{\mu\,\eta_0}.
\end{equation}
Она интегрировалась точно (с $\sin\varphi$, без замены $\cos\varphi\to1$) методом Дормана--Принса с относительной точностью $10^{-12}$; параметры условные, оба знака $\eta_0$. Среднее $\langle\delta\rangle$ и частота $\Omega$ измерялись за синхротронный период и сравнивались с~\eqref{eq:sol-deq} и~\eqref{eq:sol-omega-sum} при измеренной амплитуде первой гармоники. Во всех случаях $\langle\tau\rangle$ равно нулю с точностью $10^{-17}$.

\begin{table}[htb]
\caption{Полная система~\eqref{eq:num-sys}: счёт и формулы. $\delta_m\approx10^{-2}$.}
\label{tab:num-sol}
\small
\begin{ruledtabular}
\begin{tabular}{rrrrcccc}
$\eta_0$ & $\eta_1$ & $\Lambda$ & $\mu$ & $\langle\delta\rangle$, счёт & $\langle\delta\rangle$, \eqref{eq:sol-deq} & $\Omega/\omega_s-1$, счёт & \eqref{eq:sol-omega-sum} \\
\hline
$-0.5$ & $0$ & $0$ & $10$ & $<10^{-17}$ & $0$ & $-1.5628\times10^{-4}$ & $-1.5624\times10^{-4}$ \\
$0.4$ & $0$ & $0$ & $10$ & $<10^{-17}$ & $0$ & $-1.0001\times10^{-4}$ & $-0.9999\times10^{-4}$ \\
$-0.5$ & $0.5$ & $0$ & $1$ & $5.0002\times10^{-5}$ & $5.0000\times10^{-5}$ & $-4.3232\times10^{-5}$ & $-4.3229\times10^{-5}$ \\
$0.4$ & $-0.8$ & $0$ & $1$ & $1.0002\times10^{-4}$ & $1.0000\times10^{-4}$ & $-1.6771\times10^{-4}$ & $-1.6766\times10^{-4}$ \\
$-0.5$ & $0.5$ & $2\times10^{-5}$ & $1$ & $9.0008\times10^{-5}$ & $9.0000\times10^{-5}$ & $-8.3239\times10^{-5}$ & $-8.3229\times10^{-5}$ \\
$0.4$ & $-0.5$ & $2\times10^{-5}$ & $1$ & $1.2500\times10^{-5}$ & $1.2500\times10^{-5}$ & $-3.6043\times10^{-6}$ & $-3.6037\times10^{-6}$ \\
$-0.5$ & $0.5$ & $0$ & $10$ & $4.9999\times10^{-5}$ & $4.9996\times10^{-5}$ & $-1.9793\times10^{-4}$ & $-1.9790\times10^{-4}$ \\
$-0.5$ & $0$ & $2\times10^{-4}$ & $10$ & $4.0000\times10^{-4}$ & $4.0000\times10^{-4}$ & $-1.5628\times10^{-4}$ & $-1.5624\times10^{-4}$ \\
\end{tabular}
\end{ruledtabular}
\end{table}

Первые две строки --- только поправка фазы; частота совпадает с точной частотой маятника до $10^{-15}$, остаток относительно $-\varphi_m^2/16$ есть член $\varphi_m^4$. Остаток среднего в строках 3, 4, 7 совпадает с членом $\delta_m^4$ из~\eqref{eq:sol-a4}. В строке 6 перекрёстный член ($+6.25\times10^{-5}$) и член $5/12$ ($-6.51\times10^{-5}$) почти сокращаются, и счёт воспроизводит разность.

\subsection{Гладкие электростатические кольца}

Релятивистское движение заряда в поле $E\propto r^{-1}$ и $E\propto R^{-2}$ ($c=m=\rho=1$) интегрировалось при $\gamma_s=1.25$ и $2$. Коэффициенты $\eta_0$, $\eta_1$, извлечённые из периодов круговых орбит, совпадают с таблицей~\ref{tab:es-smooth} с точностью $10^{-8}$ (например, $\eta_1=21.75$ для сферического дефлектора при $\gamma_s=2$). Бетатронный вклад определялся по разности периодов для амплитуд $a$ и $2a$, что исключает краевую ошибку неполного бетатронного периода: $\Lambda/X=0.002$ и $0.001$ для цилиндрического дефлектора (ожидается $0$) и $-0.362$, $-0.753$ для сферического (ожидается $-\beta_s^2=-0.36$, $-0.75$).


\bibliography{references}

\end{document}
